\labsection{6}{Reinforcement learning}{sec:lab06-reinforcement}

\begin{labproject}{Learn a policy by interacting, then save it}
\textbf{Goal.} Train tabular Q-learning, evaluate the greedy policy, visualize values/actions, and save the Q-table.

\medskip\noindent\textbf{Setup.}\par
\begin{lstlisting}[style=mlpython]
from pathlib import Path

import matplotlib.pyplot as plt
import numpy as np

PROJECT_ROOT = Path.cwd().parent
MODELS = PROJECT_ROOT / "models"
MODELS.mkdir(exist_ok=True)
\end{lstlisting}

\medskip\noindent\textbf{Part A --- Tabular Q-learning in a 4$\times$4 GridWorld.}\par

Use a $4\times4$ grid: start top-left, goal bottom-right, and a trap. Ordinary moves cost $-0.04$, the goal pays $+1$, and the trap pays $-1$ and terminates. Walls leave position unchanged. \texttt{step} returns state, reward, and termination.

\begin{lstlisting}[style=mlpython]
rng = np.random.default_rng(42)

ROWS, COLS = 4, 4
START, GOAL, TRAP = (0, 0), (3, 3), (1, 3)
ACTIONS = [(-1, 0), (1, 0), (0, -1), (0, 1)]        # up, down, left, right
ACTION_NAMES = ["up", "down", "left", "right"]
ARROWS = {"up": (0, 1), "down": (0, -1), "left": (-1, 0), "right": (1, 0)}

def step(state, action):
    """Apply one action. Returns (next_state, reward, done)."""
    r, c = state
    dr, dc = ACTIONS[action]
    next_state = (min(max(r + dr, 0), ROWS - 1), min(max(c + dc, 0), COLS - 1))
    if next_state == GOAL:
        return next_state, 1.0, True
    if next_state == TRAP:
        return next_state, -1.0, True
    return next_state, -0.04, False
\end{lstlisting}

Initialize state--action Q-values to zero. Update toward reward plus discounted best next-state value, with step size $\alpha$ and discount $\gamma$. Decay exploration from $\varepsilon=1$.

\begin{lstlisting}[style=mlpython]
Q = np.zeros((ROWS, COLS, len(ACTIONS)))
alpha, gamma = 0.15, 0.95
epsilon, min_epsilon = 1.0, 0.05

def choose_action(state, explore=True):
    if explore and rng.random() < epsilon:
        return int(rng.integers(len(ACTIONS)))       # explore: random move
    return int(np.argmax(Q[state]))                  # exploit: best known move

training_returns = []
for episode in range(5000):
    state, episode_return = START, 0.0
    for _ in range(100):
        action = choose_action(state)
        next_state, reward, done = step(state, action)
        episode_return += reward
        target = reward if done else reward + gamma * np.max(Q[next_state])
        Q[state][action] += alpha * (target - Q[state][action])
        state = next_state
        if done:
            break
    training_returns.append(episode_return)
    epsilon = max(min_epsilon, epsilon * 0.9985)
\end{lstlisting}

Plot returns smoothed over 100 episodes as exploration decays.

\begin{lstlisting}[style=mlpython]
window = 100
moving = np.convolve(training_returns, np.ones(window) / window, mode="valid")
plt.figure(figsize=(6, 3.6))
plt.plot(np.arange(window - 1, len(training_returns)), moving)
plt.axhline(0.8, linestyle="--", color="grey", label="best possible return (6 moves)")
plt.xlabel("Episode")
plt.ylabel("Mean return over last 100 episodes")
plt.title("Q-learning training curve")
plt.legend()
plt.tight_layout()
plt.show()
\end{lstlisting}

Evaluate separately with exploration off, choosing the largest Q-value in each cell.

\begin{lstlisting}[style=mlpython]
successes, returns = 0, []
for _ in range(200):
    state, total = START, 0.0
    for _ in range(100):
        state, reward, done = step(state, choose_action(state, explore=False))
        total += reward
        if done:
            successes += int(state == GOAL)
            break
    returns.append(total)

print("success rate:", successes / 200)
print("mean return:", round(float(np.mean(returns)), 3))
print("learned action at start:", ACTION_NAMES[choose_action(START, explore=False)])
\end{lstlisting}

\textbf{Inspect.} Six moves yield $5(-0.04)+1=0.80$. A 100\% evaluated success rate at return 0.8 indicates the shortest route.

\medskip\noindent\textbf{Part B --- The Q-table is the model: save it and read the policy.}\par

Save the Q-table; deployment chooses each state’s highest-valued action without retraining.

\begin{lstlisting}[style=mlpython]
q_path = MODELS / "gridworld_q_table.npy"
np.save(q_path, Q)

Q_loaded = np.load(q_path)
policy = np.full((ROWS, COLS), "", dtype=object)
for r in range(ROWS):
    for c in range(COLS):
        policy[r, c] = "GOAL" if (r, c) == GOAL else "TRAP" if (r, c) == TRAP \
            else ACTION_NAMES[int(np.argmax(Q_loaded[r, c]))]
print(policy)
\end{lstlisting}

Plot greedy actions as arrows over maximum Q-values.

\begin{lstlisting}[style=mlpython]
def plot_policy(Q_table, title):
    values = Q_table.max(axis=2)
    plt.figure(figsize=(4.4, 4.2))
    plt.imshow(values, cmap="YlGn")
    plt.colorbar(label="value of the best action")
    for r in range(ROWS):
        for c in range(COLS):
            if (r, c) == GOAL:
                plt.text(c, r, "GOAL", ha="center", va="center", fontweight="bold")
            elif (r, c) == TRAP:
                plt.text(c, r, "TRAP", ha="center", va="center", color="red", fontweight="bold")
            else:
                dx, dy = ARROWS[ACTION_NAMES[int(np.argmax(Q_table[r, c]))]]
                plt.arrow(c, r, 0.3 * dx, -0.3 * dy, head_width=0.15, color="black")
    plt.xticks(range(COLS))
    plt.yticks(range(ROWS))
    plt.title(title)
    plt.tight_layout()
    plt.show()

plot_policy(Q_loaded, "Greedy policy and state values")
\end{lstlisting}

\textbf{Inspect.} Trace routes toward the goal and around the trap; compare neighboring state values.

\begin{tryit}
Change the step cost, trap, exploration schedule, or discount. Predict the policy change, retrain, and check with \texttt{plot\_policy}.
\end{tryit}
\end{labproject}
